\chapter{Resultados experimentales}\label{chap:resultados}

Este capítulo presenta los resultados obtenidos durante la evaluación del asistente institucional desarrollado. El objetivo principal es analizar el comportamiento de las distintas variantes del sistema desde tres perspectivas complementarias: la calidad de la recuperación documental, la calidad de las respuestas generadas y la utilidad funcional del orquestador.

La evaluación se organiza en varios bloques. En primer lugar, se describe brevemente el corpus utilizado, el dataset de evaluación y las métricas empleadas. A continuación, se comparan los resultados de recuperación del baseline determinista frente a la variante con recuperación híbrida y reranking. Posteriormente, se analiza la calidad de las respuestas del flujo determinista y del orquestador. Finalmente, se discuten los resultados obtenidos, destacando las mejoras observadas y las principales limitaciones del proceso de evaluación.

\section{Diseño de la evaluación}\label{sec:diseno-evaluacion}

La evaluación experimental se diseñó con el propósito de comparar el comportamiento del sistema en diferentes configuraciones. Para ello, se utilizó un corpus documental construido a partir de información pública de UNIR y un conjunto de preguntas diseñado específicamente para cubrir distintos tipos de consulta: preguntas respondibles desde un único documento, preguntas que requieren combinar información de varias fuentes, preguntas comparativas y preguntas no respondibles a partir del corpus disponible.

\subsection{Corpus y dataset de evaluación}\label{subsec:corpus-dataset}

La evaluación se realizó sobre un corpus documental construido a partir de información pública del dominio institucional. Tras el proceso de recopilación, limpieza, normalización y segmentación descrito en el capítulo anterior, el corpus final quedó formado por 26 documentos en formato Markdown. Este formato se seleccionó porque permite conservar una estructura textual sencilla, facilitar la segmentación por bloques de contenido y mantener una representación homogénea de fuentes originalmente heterogéneas.

A partir de este corpus se construyó manualmente un dataset de evaluación compuesto por 71 preguntas representativas del dominio. El objetivo de este conjunto no es constituir un benchmark general de asistentes institucionales, sino proporcionar una validación controlada del prototipo y permitir la comparación de distintas variantes del sistema bajo las mismas condiciones experimentales.

Las preguntas se diseñaron buscando cubrir escenarios habituales de uso en un contexto universitario. En concreto, se incluyeron preguntas respondibles desde un único documento, preguntas que requieren combinar información procedente de varias fuentes, preguntas comparativas y preguntas no respondibles a partir del corpus disponible. Esta última categoría se incorporó de forma deliberada para evaluar la capacidad del sistema de abstenerse cuando no existe evidencia suficiente, un comportamiento especialmente importante en asistentes institucionales donde la generación de información no fundamentada puede afectar a la confianza del usuario.

Para las preguntas respondibles se anotó manualmente la fuente esperada, es decir, el documento o documentos en los que debía encontrarse la evidencia necesaria para responder. Además, en un subconjunto de preguntas se definieron hechos esperados, entendidos como unidades mínimas de información que una respuesta correcta debería incluir. Esta anotación permitió evaluar no solo si el sistema recuperaba la fuente adecuada, sino también si la respuesta generada incorporaba los elementos informativos relevantes. A diferencia de la versión inicial del dataset, todas las preguntas respondibles del conjunto final cuentan con al menos dos hechos esperados anotados.

El tamaño del dataset, con 71 preguntas, representa una ampliación significativa respecto a versiones anteriores de trabajo. No obstante, los valores obtenidos deben entenderse como una aproximación a la evaluación del sistema y no como una medición definitiva, dado que el corpus documental cubre exclusivamente el dominio público de UNIR. Por tanto, los resultados deben interpretarse como una evaluación inicial y controlada, no como una medición definitiva del rendimiento del sistema en producción.

Con vistas a trabajos futuros, se plantea continuar ampliando el dataset de forma progresiva, priorizando preguntas de calidad sobre volumen bruto y cubriendo escenarios de consulta no representados en la versión actual. Como línea de extensión adicional, cabría ampliar el corpus documental a otra universidad de características similares, lo que permitiría validar la generalización del sistema más allá del dominio específico de UNIR.

\begin{table}[H]
\centering
\begin{tabular}{lr}
\hline
\textbf{Elemento} & \textbf{Valor} \\
\hline
Documentos finales del corpus & 26 \\
URLs semilla consideradas & 30 \\
Páginas descargadas correctamente & 30 \\
Páginas descartadas durante la limpieza & 4 \\
Formato de almacenamiento & Markdown \\
Tamaño medio de documento & 1557,92 palabras \\
Preguntas totales del dataset & 71 \\
Preguntas respondibles & 61 \\
Preguntas no respondibles & 10 \\
Preguntas con fuente esperada anotada & 61 \\
Preguntas con hechos esperados anotados & 61 \\
\hline
\end{tabular}
\caption{Resumen del corpus y del dataset utilizados en la evaluación.}
\label{tab:resumen_dataset}
\end{table}

Estos valores reflejan un proceso de recopilación acotado y controlado. Aunque el número de documentos finales no es elevado, el corpus se construyó a partir de fuentes seleccionadas manualmente y revisadas durante la limpieza, priorizando la relevancia institucional, la trazabilidad de la fuente y la utilidad para el conjunto de preguntas de evaluación. Por tanto, el objetivo no es maximizar el volumen documental, sino disponer de una base suficientemente representativa y coherente para validar el comportamiento del prototipo bajo condiciones comparables.

\begin{table}[H]
\centering
\begin{tabular}{lr}
\hline
\textbf{Tipo de pregunta} & \textbf{Número de preguntas} \\
\hline
Respondible desde un único documento & 36 \\
Respondible con combinación de documentos & 15 \\
Comparativa & 10 \\
No respondible & 10 \\
\hline
\end{tabular}
\caption{Distribución de preguntas del dataset de evaluación.}
\label{tab:distribucion_preguntas}
\end{table}

\subsection{Métricas empleadas}\label{subsec:metricas}

Para evaluar el sistema se utilizaron métricas diferentes en función del componente analizado. En la fase de recuperación documental se emplearon métricas habituales en sistemas de búsqueda de información:

\begin{itemize}
	\item \textbf{Hit@1}: indica si la fuente esperada aparece en la primera posición de los resultados recuperados.
	\item \textbf{Hit@3}: indica si la fuente esperada aparece entre los tres primeros resultados recuperados.
	\item \textbf{MRR}: mide la posición media recíproca de la primera fuente relevante recuperada.
\end{itemize}

En todas las variantes evaluadas, el número de fragmentos recuperados y proporcionados al modelo de lenguaje como contexto es $k = 5$. En la variante con recuperación híbrida y reranking, el recuperador obtiene hasta 20 candidatos iniciales (10 densos más 10 léxicos), que el reranker reduce a los 5 más relevantes antes de la generación.

Para evaluar la calidad de las respuestas se utilizaron métricas orientadas al contenido generado:

\begin{itemize}
	\item \textbf{Cobertura del \textit{ground truth}}: proporción de hechos esperados del \textit{ground truth} que aparecen cubiertos en la respuesta generada.
	\item \textbf{Tasa de abstención}: proporción de preguntas no respondibles en las que el sistema evita responder sin evidencia suficiente.
	\item \textbf{Exactitud en la detección de intención de correo}: capacidad del orquestador para identificar solicitudes que requieren activar la herramienta de generación de borradores.
\end{itemize}

\subsection{Variantes evaluadas}\label{subsec:variantes-evaluadas}

Las variantes del sistema se diseñaron como un \textit{ablation study}: cada configuración introduce un componente adicional respecto a la anterior, lo que permite aislar el efecto de cada elemento del pipeline sobre las métricas de evaluación. La Tabla~\ref{tab:variantes-evaluadas} resume las cuatro variantes propuestas.

\begin{table}[H]
\centering
\begin{tabular}{lllcc}
\toprule
\textbf{Variante} & \textbf{Recuperación} & \textbf{Reranking} & \textbf{Orquestador} & \textbf{Evaluada} \\
\midrule
A: Semántico + LLM & Densa & No & No & Sí \\
B: Híbrido + LLM & Densa + léxica & No & No & Sí \\
C: Híbrido + Reranker + LLM & Densa + léxica & Sí & No & Sí \\
D: Orquestador + Híbrido + Reranker + LLM & Densa + léxica & Sí & Sí & Sí \\
\bottomrule
\end{tabular}
\caption{Diseño del \textit{ablation study}: variantes evaluadas.}
\label{tab:variantes-evaluadas}
\end{table}

En el trabajo se evaluaron las cuatro variantes del \textit{ablation study}. Las variantes A, B y C permiten aislar el efecto individual de la recuperación híbrida y del reranker sobre la calidad de la evidencia recuperada. La variante D combina el orquestador con recuperación híbrida y reranking, y permite analizar el efecto del componente de orquestación sobre el control del flujo y la activación de herramientas.

El orquestador implementado sigue el paradigma ReAct \cite{yao2023react} (\textit{Reasoning + Acting}), en el que el modelo de lenguaje actúa como motor de decisión dentro de un bucle iterativo de \textit{Pensamiento $\to$ Acción $\to$ Observación}.
En cada paso, el LLM selecciona la herramienta a invocar —búsqueda en el corpus (\textit{buscar}), rechazo por falta de evidencia (\textit{rechazar}), redacción de correo (\textit{redactar\_correo}) o señal de respuesta suficiente (\textit{responder})— y el sistema devuelve la observación al modelo para el siguiente ciclo.
Este diseño permite que el orquestador adapte su estrategia de recuperación en tiempo de ejecución en función del contenido real de los documentos encontrados.


\section{Resultados de recuperación}\label{sec:resultados-recuperacion}

En esta sección se comparan los resultados obtenidos por las variantes A, B y C del \textit{ablation study}, evaluadas sobre las preguntas respondibles del dataset que cuentan con una fuente esperada, excluyendo las preguntas no respondibles.

\begin{table}[H]
\centering
\begin{tabular}{lcccccc}
\toprule
\textbf{Variante} & \textbf{Hit@1} & \textbf{$\Delta$} & \textbf{Hit@3} & \textbf{$\Delta$} & \textbf{MRR} & \textbf{$\Delta$} \\
\midrule
A: Semántico + LLM & 0.574 & -- & 0.852 & -- & 0.712 & -- \\
B: Híbrido (RRF) + LLM & 0.525 & $-$0.049 & 0.902 & +0.050 & 0.706 & $-$0.006 \\
C: Híbrido (RRF) + Reranker + LLM & 0.557 & $-$0.017 & 0.918 & +0.066 & 0.738 & +0.026 \\
\bottomrule
\end{tabular}
\caption{Resultados de recuperación del \textit{ablation study}.}
\label{tab:resultados-recuperacion}
\end{table}

Los resultados muestran un patrón diferenciado según el componente incorporado. La variante A (semántico puro) obtiene Hit@1 = 0.574, Hit@3 = 0.852 y MRR = 0.712 como línea base, con buena precisión en la primera posición pero cobertura moderada en el top-3.

La variante B (recuperación híbrida con \textit{Reciprocal Rank Fusion}) obtiene Hit@1 = 0.525 ($-$0.049 respecto a A), Hit@3 = 0.902 (+0.050) y MRR = 0.706 ($-$0.006). La combinación densa y léxica introduce mayor diversidad de candidatos, lo que incrementa la presencia de documentos relevantes en el top-3 (+5,0 puntos porcentuales), pero puede desplazar documentos altamente relevantes de la primera posición cuando BM25 introduce candidatos léxicamente similares con menor relevancia semántica.

La variante C (híbrido RRF + reranker \textit{BAAI/bge-reranker-v2-m3}) obtiene Hit@1 = 0.557 ($-$0.017 respecto a A), Hit@3 = 0.918 (+0.066) y MRR = 0.738 (+0.026). El reranker multilingüe recupera parcialmente la precisión en primera posición y eleva la MRR por encima de la línea base, lo que indica que el modelo \textit{cross-encoder} es capaz de reordenar correctamente los candidatos del pool híbrido. Hit@3 alcanza el valor más alto de las tres variantes (0.918), confirmando que la combinación de recuperación híbrida con reranking proporciona el pool de evidencia más rico para la generación.

\section{Resultados de calidad de respuesta}\label{sec:resultados-respuesta}

Además de evaluar la recuperación documental, se analizaron las respuestas generadas por el sistema. Esta evaluación permite observar si las mejoras introducidas en la arquitectura se traducen en respuestas más completas, más seguras o más controladas.

La Tabla~\ref{tab:resultados-calidad-respuesta} compara el comportamiento del flujo determinista y del orquestador en tres dimensiones: cobertura del \textit{ground truth}, abstención ante preguntas no respondibles y detección de solicitudes de correo.

\begin{table}[H]
\centering
\begin{tabular}{lcc}
\toprule
\textbf{Métrica} & \textbf{A: Determinista} & \textbf{C: Orquestador} \\
\midrule
Cobertura del \textit{ground truth} & 0.438 & 0.312 \\
Tasa de abstención en preguntas no respondibles & 0.000 & 1.000 \\
Exactitud en detección de intención de correo & 0.000 & 0.833 \\
\bottomrule
\end{tabular}
\caption{Comparativa de calidad de respuesta entre el flujo determinista y el orquestador.}
\label{tab:resultados-calidad-respuesta}
\end{table}

Los resultados muestran un comportamiento diferenciado entre ambas variantes. El flujo determinista obtiene una mayor cobertura del \textit{ground truth} en preguntas respondibles, con un valor de 0.438 frente al 0.312 obtenido por el orquestador. Este resultado indica que, en esta configuración, el flujo determinista tiende a producir respuestas más alineadas con los hechos esperados definidos en el dataset.

Sin embargo, el orquestador mejora de forma clara el comportamiento ante preguntas no respondibles. Mientras que el flujo determinista no se abstiene en este tipo de preguntas, el orquestador alcanza una tasa de abstención de 1.000. Este resultado es especialmente relevante en un asistente institucional, ya que responder sin evidencia suficiente puede producir información incorrecta o poco fiable.

Además, el orquestador permite detectar solicitudes accionables, como la generación de borradores de correo. En esta dimensión alcanza una exactitud de 0.833, frente al 0.000 del flujo determinista, que no incorpora mecanismos de planificación ni activación de herramientas. Por tanto, aunque el orquestador no mejora la cobertura factual en esta evaluación, sí aporta capacidades relevantes de control, abstención y uso de herramientas.

\section{Evaluación funcional del orquestador}\label{sec:evaluacion-funcional-orquestador}

La evaluación del orquestador no se limita únicamente a métricas de respuesta, ya que su objetivo principal no es sustituir al recuperador, sino coordinar el flujo de ejecución del asistente. Por ello, se realizaron pruebas funcionales orientadas a comprobar si el sistema identifica correctamente el tipo de consulta y selecciona la acción adecuada.

La Tabla~\ref{tab:validacion-funcional-orquestador} resume los resultados obtenidos en los principales escenarios, expresados mediante métricas objetivas extraídas de los registros de evaluación.

\begin{table}[H]
\centering
\small
\begin{tabular}{lcccc}
\toprule
\textbf{Tipo de consulta} & \textbf{n} & \textbf{Enrutamiento correcto} & \textbf{Llamadas herramienta (media)} & \textbf{Abstención} \\
\midrule
Un documento & 15 & 13/15 (0,87) & 1,0 & 0/15 (0,00) \\
Multi-documento & 6 & 2/6 (0,33) & 1,3 & 0/6 (0,00) \\
Comparativa & 4 & 4/4 (1,00) & 1,8 & 0/4 (0,00) \\
No respondible & 5 & \textemdash & \textemdash & 5/5 (1,00) \\
Solicitud de correo & 6 & 5/6 (0,83) & \textemdash & \textemdash \\
\bottomrule
\end{tabular}
\caption{Evaluación funcional del orquestador con métricas objetivas ($n$ = número de preguntas de cada tipo).}
\label{tab:validacion-funcional-orquestador}
\end{table}

Los resultados muestran un comportamiento heterogéneo según el tipo de consulta. Las preguntas comparativas alcanzan el enrutamiento perfecto (4/4), y las de un único documento obtienen un 87\,\% de enrutamiento correcto con una media de una llamada a herramienta por consulta. Las preguntas no respondibles son gestionadas correctamente en todos los casos, con una tasa de abstención del 100\,\%. Las solicitudes de correo se detectan en 5 de 6 casos (83\,\%). Sin embargo, el enrutamiento de preguntas multi-documento resulta el más problemático, con solo un 33\,\% de clasificaciones correctas, lo que sugiere que este tipo de consulta resulta especialmente difícil de identificar de forma automática y requeriría mayor esfuerzo de ajuste en iteraciones futuras.

\section{Discusión de resultados}\label{sec:discusion-resultados}

Los resultados obtenidos permiten extraer varias conclusiones relevantes sobre el comportamiento del sistema.

En primer lugar, la incorporación de recuperación híbrida y reranking mejora el rendimiento del sistema en las métricas de recuperación. La variante B (híbrido RRF) muestra que la fusión léxico-semántica amplía la cobertura en el top-3: Hit@3 pasa de 0.852 a 0.902 (+0.050) respecto a A, aunque Hit@1 desciende ligeramente ($-$0.049) y MRR se mantiene prácticamente igual ($-$0.006). Este patrón indica que BM25 introduce candidatos léxicamente relevantes que mejoran la presencia de la fuente esperada en posiciones intermedias, pero no siempre en primera posición. La variante C (híbrido RRF + reranker \textit{BAAI/bge-reranker-v2-m3}) obtiene Hit@1 = 0.557 ($-$0.017 respecto a A), Hit@3 = 0.918 (+0.066) y MRR = 0.738 (+0.026). El reranker multilingüe reordena el pool híbrido de forma más precisa: Hit@3 alcanza su valor máximo y MRR supera al baseline, confirmando que el modelo \textit{cross-encoder} aporta una discriminación adicional que no ofrece ni la búsqueda semántica ni la léxica por separado.

En segundo lugar, el orquestador no mejora la cobertura del \textit{ground truth} respecto al flujo determinista en la configuración evaluada. Este resultado puede explicarse por varios factores. Por un lado, el orquestador adopta un comportamiento más conservador, especialmente cuando la evidencia recuperada no es suficientemente clara. Por otro lado, la métrica de cobertura del \textit{ground truth} depende de la coincidencia entre la respuesta generada y un conjunto concreto de hechos definidos previamente, por lo que puede penalizar respuestas correctas formuladas de manera más general o prudente. Además, el modelo generativo empleado es de tamaño reducido y los prompts del orquestador no han sido objeto de un ajuste sistemático, lo que puede introducir variabilidad en la cobertura factual con independencia del componente de recuperación. Ajustar los prompts y explorar modelos generativos de mayor capacidad son líneas de mejora directas para incrementar esta métrica en iteraciones futuras.

En tercer lugar, el orquestador sí aporta mejoras claras en dimensiones que no quedan reflejadas únicamente mediante métricas clásicas de recuperación. Su capacidad para abstenerse ante preguntas no respondibles reduce el riesgo de generar respuestas no fundamentadas. Asimismo, la detección de solicitudes de correo demuestra que la arquitectura puede extenderse hacia funcionalidades accionables, manteniendo un flujo controlado y trazable.

En conjunto, los resultados sugieren que las mejoras sobre el recuperador tienen un impacto directo en la calidad de la evidencia recuperada, mientras que el orquestador aporta valor principalmente en términos de control del flujo, seguridad de la respuesta y extensibilidad funcional. Por tanto, ambas líneas de mejora son complementarias: el reranking contribuye a recuperar mejor información, mientras que el orquestador permite decidir cómo utilizarla.

Por tanto, los resultados deben interpretarse desde la contribución global de la arquitectura: la mejora del retrieval incrementa la calidad de la evidencia disponible, mientras que la orquestación permite controlar cómo se utiliza dicha evidencia y cuándo debe activarse una capacidad accionable. La aportación principal no reside en un único componente aislado, sino en la integración modular y trazable de estos elementos dentro de un asistente institucional supervisado.

\section{Limitaciones de la evaluación}\label{sec:limitaciones-evaluacion}

La evaluación realizada presenta algunas limitaciones que deben tenerse en cuenta al interpretar los resultados. En primer lugar, el tamaño del corpus y del dataset es reducido, por lo que los valores obtenidos deben entenderse como una aproximación inicial al comportamiento del sistema y no como una medición definitiva.

En segundo lugar, la evaluación de calidad de respuesta se basa en un subconjunto de preguntas con hechos esperados definidos manualmente. Este enfoque permite realizar una comparación semiautomatizada, pero puede penalizar respuestas que sean parcialmente correctas o que expresen la información con una formulación distinta a la esperada.

En tercer lugar, el modelo generativo utilizado en la evaluación es ligero, lo que puede afectar a la calidad final de las respuestas. En consecuencia, parte de los errores observados pueden deberse no solo al diseño del pipeline RAG, sino también a las limitaciones del modelo de lenguaje empleado.

Finalmente, la evaluación funcional del orquestador permite comprobar que los flujos principales se comportan como se esperaba, pero sería conveniente ampliarla en trabajos futuros con un conjunto mayor de consultas y métricas específicas de routing, planificación y uso de herramientas.

\section{Síntesis de resultados}\label{sec:sintesis-resultados}

A modo de síntesis, los experimentos realizados muestran que la variante C (híbrido RRF + reranker) es la mejor configuración global para recuperación: obtiene el Hit@3 más alto (0.918, +0.066 respecto a A) y la MRR más alta (0.738, +0.026 respecto a A). El \textit{ablation study} revela que ambos componentes contribuyen de forma diferenciada y complementaria: la variante B (híbrido RRF sin reranker) mejora principalmente la cobertura en el top-3 (Hit@3 +0.050), mientras que la variante C añade el reranker multilingüe que no solo amplía la cobertura sino que también mejora la calidad del ranking (MRR +0.026 sobre A). Hit@1 desciende ligeramente en B y C respecto a A, lo que sugiere que el pool de candidatos ampliado por BM25 puede introducir ruido en la primera posición que el reranker no siempre compensa completamente.

Por otro lado, la comparación entre el flujo determinista y el orquestador muestra que el segundo no incrementa la cobertura del \textit{ground truth} en esta configuración, pero sí mejora de forma notable la abstención ante preguntas no respondibles y permite incorporar capacidades accionables. Estos resultados respaldan la utilidad de una arquitectura modular, en la que las mejoras de recuperación y las capacidades de orquestación cumplen funciones distintas pero complementarias dentro del asistente institucional.